% A product chain is one n-ary product, a division chain left-associative, and an explicit MulOp takes one
% factor on its right (Perl's `moreFactors`, MathGrammar:252-258: `MulOp Factor` and `Factor`, each `ApplyNary`d): a
% juxtaposition after an n-ary MulOp chain multiplies the whole chain, `a\cdot b\cdot c d` (a·b·c)·d and `a/b/c d`
% ((a/b)/c)·d, as the binary `a\cdot b c` (a·b)·c — SYNC_STATUS "Math-parse residuals" (14), 57cj.19; before it the n-ary tuck in
% `infix_apply_nary` took the juxtaposed product whole, a·b·(c d), a/b/(c d).
\documentclass{article}
\usepackage{amsmath,amssymb,bm,stmaryrd,xcolor}
\begin{document}

\section{The class}
$a\cdot b c$ $a\cdot b\cdot c d$ $a\cdot b\cdot c\cdot d e$ $a\cdot b\cdot c d e$ $a b\cdot c\cdot d e$
$a\cdot b\cdot c d\cdot e f$ $a\cdot b\cdot c\,d$ $a\cdot b\cdot c d+e$ $x=a\cdot b\cdot c d$ $-a\cdot b\cdot c d$
$a\times b\times c d$ $a\star b\star c d$ $a\otimes_k b\otimes_k c d$ $a\cdot b\times c d$ $a/b\cdot c d$
$a/b c$ $a/b/c d$ $a/b/c d/e$ $a\cdot b\cdot(c d)$ $a/(b c)$

% (Q11, user ruling 2026-10-01: juxtaposition binds tighter than a large MULOP) a large product operator — ⊗, ⊙ and the
% circled and boxed family ⊘ ⊛ ⊚ ⊠ ⊡ — keeps its juxtaposed operand whole, as a BINOP does (#393, divergence #396):
% `2\Lambda_1\otimes\Lambda_1\otimes 2\Lambda_1\otimes\Lambda_1` (2Λ₁)⊗Λ₁⊗(2Λ₁)⊗Λ₁ (2605.17901), `g\otimes w\otimes\sigma'(\gamma_i)`
% (2605.01702), `P\odot P\odot P_\theta(x|y)` (2605.00423), `c\boxtimes T c` (2605.29990), a decorated one too
% (`R\otimes_{\mathbb C}\mathbb C G`, 2605.14864), and in a bare argument (`\sin x\otimes y z`); before an integral's
% differentials it takes the integrand's factors before them (`\int f\otimes g h\,dx` ∫((f⊗(g h))·dx)). `\cdot`, `\times`,
% `\star`, `\circ` keep Perl's one factor ((f∘g)·x); the semidirect products, ∐ and the circles join the class (below).
\section{A large product operator keeps its juxtaposed operand}
$2\Lambda_1\otimes\Lambda_1\otimes 2\Lambda_1\otimes\Lambda_1$ $g\otimes w\otimes\sigma'(\gamma_i)$ $a\otimes 2b$ $P\odot P\odot P_\theta(x|y)$
$a\otimes b c$ $a\otimes b\otimes c d$ $x\odot y z$ $A\boxtimes B C$ $a\oslash b c$ $a\circledast b c$ $a\boxdot b c$ $f\circledcirc g h$
$c\boxtimes T c$ $R\otimes_{\mathbb C}\mathbb C G$ $a\otimes b c+d$ $x=a\otimes b c$ $\partial_x u\otimes v w$ $\sum_i a_i\otimes b_i c$
$\sin x\otimes y z$ $\log a\otimes b c$ $\nabla u\otimes v w$ $\sin\log x\otimes y z$
$\int f\otimes g\,dx$ $\int f\otimes g h\,dx$ $\int f\otimes g h\,dx\,dy$ $\int_{s<u_1<u_2<t}dX_{u_1}\otimes dX_{u_2}$ $\int_X F(x)\otimes d\nu(x)$
% (57cj.20.Q11 review) the integrand closes before the letter `d` too, so a definite integral keeps its differential
% (`\int_0^1 f\otimes g h\,dx` read ∫(f⊗(g·h·d·x))); a trig argument's chain goes on in the operand (`\sin x\otimes 2\pi t`
% read sin@(x⊗(2π))·t), and a run of ellipses between two items joins it; ⦸ is of the family. Witnesses: the
% stochastic-integral rows 2605.20541 (`dX_{u_1}\otimes\cdots\otimes dX_{u_a}`), 2605.25146 (`F(x)\otimes d\nu(x)`)
$\int_0^1 f\otimes g h\,dx$ $\int_a^b f\odot g h\,dx\,dy$ $\int f\otimes g h\,dx\,dy\,dz\,dw$ $\int_0^1 f\otimes g h\,d^3x$
$\int_0^1 f\boxtimes g h\,dx$ $\int_0^1 f\otimes g h\,dx\,k$ $\sin x\otimes 2\pi t$ $\cos\theta\otimes 2\phi\psi$ $\tan x\otimes y3z$
$\log x\otimes y\cdots z$ $\sin x\otimes y\cdots z$ $x\otimes y\cdots z$ $\log x\otimes y\cdots$ $a\varobslash b c$
% controls: an application keeps its group, an operand on the left is the juxtaposition already, and the other MULOPs
% keep their one factor
$A\otimes B(x)$ $\nabla u\otimes\nabla v$ $a b\otimes c$ $a\cdot b c$ $f\circ g x$ $a\times b c$ $a\star b c$
% witness: the Q11 scope ruling (user, 2026-10-01): the semidirect products ⋉ ⋊ ⋋ ⋌, the coproduct ∐ (`\amalg`) and the circles
%          ○ ◯ join the large MULOPs (#396; mathabx's □ and `\pluscirc` in large_mulops_mathabx.tex) — 2605.11552
%          `\pi_1(M_0)=\pi_1(T^3)\rtimes\pi_1(N_0)`, 2605.15276 `\Gamma_{\rm mirror}\ltimes{\rm H}^{2h^{1,1}+3}(\mathbb Z)`; controls (no juxtaposed operand, unchanged): 2605.27086
%          `\operatorname{Gau}(TM)\rtimes\operatorname{Diff}(M)`, 2605.12221 `B\rtimes_r^{(\alpha,\sigma)}W` (its row appends an `x`)
% engines: rust m45: (G⋉H)·x, ((π₁T³)⋊π₁)·N₀ (Perl's one factor)
% expect:  0 errors; each keeps its juxtaposed operand whole, before an integral's differentials the integrand's factors
\section{The semidirect products, the coproduct and the circles}
$G\ltimes H x$ $G\rtimes H x$ $a\leftthreetimes b c$ $a\rightthreetimes b c$ $A\amalg B C$ $f\bigcirc g h$ $f\varbigcirc g h$
$\pi_1(M_0)=\pi_1(T^3)\rtimes\pi_1(N_0)$ $B\rtimes_{r}^{(\alpha,\sigma)}W x$ $\operatorname{Gau}(TM)\rtimes\operatorname{Diff}(M)$
$\Gamma_{\rm h}=\mathsf{SL}(2,\mathbb{Z})\ltimes\left(\Gamma_{\rm mirror}\ltimes{\rm H}^{2h^{1,1}+3}(\mathbb{Z})\right)$ $\int f\ltimes g h\,dx$ $G\ltimes 2H$
$\sin x\ltimes y z$ $\log x\rtimes y\cdots z$ $\log x\rtimes y\cdots$ $G\rtimes_\alpha H x$ $A\amalg B\amalg C D$

\section{Scripts, fences, fractions and postfix operators}
$a\cdot b\cdot c^2 d_1$ $a\cdot b\cdot\frac{c}{e}d$ $a\cdot b\cdot c!d$ $a\cdot b\cdot c'd$ $a\cdot b\cdot|c|d$
$a\cdot b\cdot f(x)g(y)$ $a\cdot b\cdot c\mathbb{E}[X]$ $S_{w_1}\circ\cdots\circ S_{w_n}(0)$ $a\cdot b\cdot c d\cdots$

\section{Inside operands}
$\sum_i a_i\cdot b_i\cdot c_i d_i$ $\int a\cdot b\cdot c\,dx$ $\partial_x u\cdot v\cdot w z$ $\sin a\cdot b\cdot c d$
$\nabla u\cdot v\cdot w z$ $\log x\cdot y\cdot z w$

% (57cj.19 reviews) A BINOP keeps its juxtaposed operand whole, where Perl's `MulOp : BINOP` (MathGrammar:688) takes one
% factor (divergence #393: a `\mathbin` of unknown meaning is no product — `KX\mathbin{\|}(I-K)X` 2605.31129,
% `[WX_i\mathbin{\|}WX_j\mathbin{\|}w_{ij}]` 2605.31315, 2605.08689; Q11 extends it to the large MULOPs, above); before an integral's
% differentials (a `d`'s, not a ∂ operator's) after an integrand it takes the integrand's factors before them, as a large MULOP
% does (user ruling 2026-10-01, revisiting 57cj.19.3-.8's one factor: `\int f\boxast g h\,dx` ∫((f⧆(g h))·dx), was
% ∫((f⧆g)·h·dx); `\int f\boxast g\,dx` ∫((f⧆g)·dx); an operand that opens with a differential stays whole,
% `\int f\boxast dx\,g` ∫(⧆(f, dx·g)), where Perl splits); around a big operator it is a
% MulOp, as Perl (`a\boxast\sum_i b_i c_i` ⧆(a, ∑…), was unparsed; `\sum_i a_i\boxast b_i` ∑(a_i⧆b_i), was ⧆(∑a_i, b_i));
% only the same operator flattens (Perl `isSameExpr`: meaning, value and mathstyle; `a/b\div c` was one n-ary `/` that
% dropped the ÷); a coloured or `\mathbin` MULOP flattens as the bare one does.
\section{A BINOP, a different operator of the same meaning, a coloured operator}
$a\boxast b c$ $a\boxast b\boxast c d$ $a\mathbin{\#}b\mathbin{\#}c d$ $a\boxast_k b c$ $a\boxast 2b$ $a/b\div c$
$a\div b/c d$ $a{\color{red}\cdot}b{\color{red}\cdot}c d$ $a\mathbin{\cdot}b\mathbin{\cdot}c d$
$\sum_i a_i\boxast b_i$ $\int_0^1 f\boxast g\,dx$ $\int f\mathbin{\|}g\,dx$ $\int f\boxast g\,dx\,dy$ $\int f\boxast g h\,dx$
$\int_X f\boxast g\,d\mu(x)$ $\int f\boxast g\,dx\,(1+h)$ $K X\mathbin{\|}Y\partial_t u$ $a\boxast b\,dx$
$\int f\boxast g\,\partial_t u$ $\int f\boxast dx\,dy$ $\int f\boxast dx\,g$ $\int dx\,f\boxast dy\,g$ $\int f\boxast dx\,g\,dy$
$\int dx\,f\boxast g\,dy$
$\int f\boxast g\boxast h k\,dx$ $\int f\boxast g\mathbin{\#}h k\,dx$ $\int f\otimes g\otimes h k\,dx$
$\int f\mathbin{\#}g h\,dx\,dy$ $\int_0^1 f\boxast 2g\,dx$ $\int f\boxast g h\,d^3x$ $\int f\boxast g h\,dx\,k$ $\int_X f\boxast g h\,d\mu(x)$
$a\boxast\sum_i b_i c_i$ $\sin x\boxast\cdots$
$K X\mathbin{\|}(I-K)X$ $[W X_i\mathbin{\|}W X_j\mathbin{\|}w_{ij}]$

\section{A letter applied to a group ends the chain}
$a_1\cdot b\cdot S(a_2)$ $a\cdot S(a_2)$ $\Gamma\cdot L\cdot M\cdot r(d+l)$

% witness: the 57cj.18 review's mine of the 3,003 A/B sources — 2605.00321, 2605.02211, 2605.02221, 2605.03798, 2605.05322, 2605.05390, 2605.07323, 2605.08833, 2605.08915, 2605.10640, 2605.10723, 2605.11288, 2605.13025, 2605.13511, 2605.14864, 2605.15530, 2605.16690, 2605.16804, 2605.18055, 2605.19033, 2605.19289, 2605.23084, 2605.24688, 2605.24740, 2605.24758, 2605.24816, 2605.25789, 2605.26484, 2605.29453, 2605.31500
\section{Product chains from the A/B}
$\frac{1}{N(1-p)} \cdot N \cdot P(m_i=0)$
$\lambda \cdot \beta \cdot \hat{\delta}(\mathbf{p})$
$a_{ij}\cdot 1\cdot e_ie_j$
$\frac{1}{1 + 66 \epsilon'} \cdot (1 - \epsilon ) \cdot \mathrm{OPT}(G)$
$(a_{1}\star x_{1})\cdot x_{2}\cdot(a_{2}\star y)\cdot S(x_{3})$
$(a_{1}\rightharpoonup x_{1})\cdot(a_{2}\rightharpoonup y_{1})\cdot S(y_{2})$
$a_{1}\cdot b\cdot S(a_{2})$
$a_{1}\cdot(x\star b)\cdot S(a_{2})$
$a_{1}\cdot(a_{2}\rightharpoonup x)\cdot S(a_{3})$
$k_1\cdot (h\star i)\cdot S(k_2)$
$h_1\cdot b\cdot S(h_2)$
$x_{1}\cdot(i\star h)\cdot S(x_{2})$
$1\cdot (1\bullet a_3)\cdot S(a_4)$
$a_1\cdot b_1\cdot S(a_2)$
$(x_{1}\rightharpoonup y_{1})\cdot(x_{3}\rightharpoonup a_{1})\cdot S(x_{2}\rightharpoonup y_{2})$
$x_1\cdot (a\star y)\cdot S(x_2)$
$x_1\cdot (y\star a)\cdot S(x_2)$
$\mathbf{T}_{L_T \leftarrow W} \cdot \mathbf{T}^t_{W \leftarrow C_k} \cdot \pi^{-1}(\mathbf{p}_j^t)$
$2^{13}\cdot 3^7\cdot 5^3\cdot 7^2\cdot 337$
$0.5 \cdot 0.813984 \cdot x_0^2 \cdot 0.813984$
$\gamma_n \cdot (1-\alpha) \cdot 2^{\alpha-1} \cdot h_n^{(-\alpha, 0)} \delta_{mn}$
$(t\xi) \cdot \xi \cdot P_n^{(-\alpha, 0)}(2\xi - 1)$
$l \cdot \boldsymbol K \cdot \Delta \boldsymbol u$
$s \cdot l \cdot \boldsymbol{m}_\theta$
$z_s\cdot y_{o,s}\cdot\delta_s(t)$
$m_i \cdot d_i \cdot Q(o_i)$
$n^m \cdot n^{(1-\zeta)\ell}\cdot 2^\ell n^{(6\epsilon-1)\ell}$
$\eta^2 \cdot e \cdot \xi_h^2(s)$
$\frac{1}{2}\cdot 60\cdot 80$
$\frac{1}{2}\cdot 100\cdot XW$
$E^{-1} \cdot E \cdot Sa$
$\frac{8(4+\lambda_\Phi)L^2}{\lambda_\Phi^2}\cdot\frac{\lambda_\Phi\alpha_\ell}{4\lambda\beta_\ell} \cdot e_\ell$
$\Gamma \cdot \mathcal{B} \cdot L(M \cdot d + K_c \cdot d \cdot l)$
$\Gamma \cdot L \cdot M \cdot r(d+l)$
$C \cdot L \cdot M \cdot r(d+l)$
$L \cdot p \cdot R(d+l)$
$L \cdot M \cdot r(d+l)$
$2 \cdot L \cdot M \cdot r(d+l)$
$C \cdot L \cdot R^2(d + l + R)$
$C \cdot L \cdot M \cdot r(d + l)$
$\Gamma \cdot \mathcal{B} \cdot L (M \cdot d + K_c \cdot d \cdot l)$
$\Gamma \cdot L \cdot M \cdot r(d + l)$
$H_{ij}^k \cdot \tilde v_{ij}^k \cdot I(z_{ij}^k > 0)$
$H_{ij}^k \cdot \tilde{v}_{ij}^k \cdot I(z_{ij}^k > 0)$
$0.5 \cdot (d_1 + d_2) \cdot \Delta t$
$2 \cdot \frac{N-1}{N^2} \cdot (N-1) \cdot \frac{N-2}{N-1} \sigma^2$
$(\bm{u}) \cdot \mathbf{K} \cdot \mathrm{diag}(\bm{v})$
$2^3\cdot 1229\cdot 20225677$
$(z, \kappa_i, \tau)\cdot \Omega \cdot b_{IT}(z)$
$i \cdot i \cdot P(s, a_1, s'')$
$\mathbf{v} \cdot (\bm{\sigma} - \rho \mathbf{v} \mathbf{v}) \cdot \mathbf{n} d$
$N_t \cdot L \cdot l^2d$
$\frac{\ln T}{\mathrm{KL}(\nu_i,\nu_1)}\cdot \frac{c^2}{(c-2)^2}\cdot \rho_{i,\min}(\frac{\Delta_i}{c-2})$
$v \cdot 0 \cdot \mathbf{0}^\top$
$d_{\max} \cdot d_{\max} \cdot L^2 B_x^2$
$M \cdot D \cdot si$
% witness: 2605.14476 (`12/z/ma`, `12/c/pz`), 2605.04903 (`828/793/733 f`)
\section{Division chains from the A/B}
$828/793/733 f$
$L_w/d^2/10^{37}$
$12/z/ma$
$12/c/pz$
$80/10/10$ $107/89/77/60/60/60/47$
\end{document}
